\hypertarget{limitations-and-scope}{%
\section{Limitations and scope}}

Both sources are selected historical archives. Weibo's held-out root users come from one publication day; SEISMIC's chronological split operates within a popularity-selected release. The two sources strengthen the empirical scope relative to a single archive, but models are trained separately and no cross-platform transfer experiment is performed. Complete collection coverage, live historical feature availability and independence of underlying events are unverified.

The six intervals are conditional on fixed fitted objects and the chosen resampling units. They do not incorporate training or tuning uncertainty. Root-user grouping addresses one identifiable dependence structure on Weibo; SEISMIC lacks comparable user grouping, and common events can connect roots in both sources. The large-count cases and duplicate sensitivity warrant particular care in extrapolating average effects.

Model grids, three fixed stochastic seeds and numerical eligibility rules define the procedures being evaluated. They do not exhaust each family's achievable performance. I1 increases input dimension and can change effective flexibility. The dimension-matched control supplies evidence against the particular uninformative augmentation tested, while leaving differences in predictor structure and effective capacity. Its intervals condition on three fixed noise realizations. The study does not separate the seven timing components or establish equivalence between noise and I0. The numerical contract excluded some validation candidates; retaining those exclusions makes the selection process explicit, but conclusions remain conditional on that contract. Marginal calibration diagnostics do not guarantee conditional calibration or operational decision value.

Both supplementary analyses were specified after the primary test results were known. Stratum intervals are exploratory pointwise summaries; the noise contrasts use their separately specified multiplicity family. They are not an independent replication on fresh test roots. The six original comparisons remain unchanged. Finally, the development history informed the prospectively frozen two-source study. Protocols were recorded before their respective performance evaluations, with metadata access and training pilots disclosed. This is not an externally registered report. The work concerns recorded growth and provides no intervention-based evidence about rumor containment. {\small\textit{Preparation:~\cite{openai2026}.}}
